\documentclass[12pt,a4paper]{article}

\usepackage[spanish]{babel}
\usepackage[utf8]{inputenc}
\usepackage[T1]{fontenc}
\usepackage{geometry}
\usepackage{amsmath}
\usepackage{amssymb}
\usepackage{booktabs}
\usepackage[ruled,vlined,linesnumbered]{algorithm2e}
\usepackage[hidelinks]{hyperref}
\usepackage{float}
\usepackage{graphicx}
\newcommand{\maybeimage}[3]{%
  \IfFileExists{#1}{\includegraphics[#2]{#1}}%
  {\fbox{\parbox{0.7\textwidth}{\centering\vspace{1cm}[Figura pendiente: #3]\vspace{1cm}}}}%
}
\usepackage{caption}
\usepackage{parskip}
\usepackage{setspace}
\usepackage{titlesec}
\usepackage{xcolor}
\usepackage{array}

\geometry{top=2.5cm, bottom=2.5cm, left=2.5cm, right=2.5cm}
\setlength{\parskip}{6pt}

% Ajuste de espaciado en secciones
\titlespacing*{\section}{0pt}{14pt}{6pt}
\titlespacing*{\subsection}{0pt}{10pt}{4pt}

% Comando para valores pendientes
\newcommand{\placeholder}[0]{\textit{---}}

% Configuración de algorithm2e
\SetKwInput{KwIn}{Entrada}
\SetKwInput{KwOut}{Salida}
\SetKw{KwReturn}{devolver}
\SetKwFor{Para}{para}{hacer}{fin para}
\SetKwFor{Mientras}{mientras}{hacer}{fin mientras}
\SetKwIF{Si}{SiNo}{OSiNo}{si}{entonces}{si no si}{si no}{fin si}
\SetKwFunction{fitness}{fitness}
\SetKwFunction{fix}{fix}
\SetKwFunction{createSolution}{crearSolución}
\SetKwFunction{shuffle}{barajar}

%-----------------------------------------------------------------------
\begin{document}

%-----------------------------------------------------------------------
% PORTADA
%-----------------------------------------------------------------------
\begin{titlepage}
  \centering
  \vspace*{2cm}
  {\Large \textbf{Universidad de Granada}\\[0.3cm]
          \textbf{ETSIIT Informática y Telecomunicación}\\[0.3cm]
          Dpto.\ de Ciencias de la Computación e Inteligencia Artificial}

  \vspace{2.5cm}
  {\LARGE \textbf{Metaheurísticas}}\\[0.5cm]
  {\Large Práctica 1}\\[0.3cm]
  {\large Técnicas de Búsqueda Local y Algoritmos Greedy\\
           para el Problema de Portfolio}

  \vspace{1.5cm}
  {\large Curso 2025--2026}

  \vfill
  \begin{tabular}{rl}
    \textbf{Nombre:}  & Jose Angel Carretero Montes \\[4pt]
    \textbf{DNI:}     & 25628892S \\[4pt]
    \textbf{Email:}   & \texttt{joseangel1@correo.ugr.es} \\[4pt]
    \textbf{Grupo:}   & Grupo 3 \\
  \end{tabular}
  \vspace{1.5cm}
\end{titlepage}

%-----------------------------------------------------------------------
% ÍNDICE
%-----------------------------------------------------------------------
\tableofcontents
\newpage

%-----------------------------------------------------------------------
\section{Descripción del Problema}
%-----------------------------------------------------------------------

El \textbf{Problema de Portfolio} consiste en determinar la asignación óptima de capital
entre $N$ activos financieros con el objetivo de maximizar el retorno esperado
minimizando simultáneamente el riesgo.

\subsection*{Representación de los datos}

Para cada activo $i \in \{1,\ldots,N\}$ y día $t \in \{1,\ldots,T\}$ se dispone
del retorno diario $\mu_{ti}$, calculado como el ratio de mejora relativa del valor:
\[
  \mu_{ti} = \frac{V_{ti} - V_{(t+1)i}}{V_{ti}}
\]
Los datos empleados cubren el periodo 2015--2024 para entrenamiento y el año 2025
para evaluación de generalización.

\subsection*{Función objetivo}

Dado un vector de pesos $w = (w_1,\ldots,w_N)$, la calidad de la cartera se mide como:
\[
  \text{fitness}(w) = \text{Beneficio}(w) - \lambda \cdot \text{Riesgo}(w)
\]
donde:
\[
  \text{Beneficio}(w) = \sum_{i=1}^{N} w_i \cdot \frac{1}{T}\sum_{t=1}^{T}\log(1+\mu_{ti})
  \qquad
  \text{Riesgo}(w) = \sum_{i=1}^{N}\sum_{j=1}^{N} w_i \cdot \Sigma_{ij} \cdot w_j
\]
siendo $\Sigma_{ij} = \frac{1}{T}\sum_{t=1}^{T}(\mu_{ti}-\bar{\mu}_i)(\mu_{tj}-\bar{\mu}_j)$
la covarianza entre activos $i$ y $j$, y $\lambda = 500$ el factor de penalización
del riesgo.

\subsection*{Restricciones}

\[
  \sum_{i=1}^{N} w_i = 1 \qquad
  w_i = 0 \;\vee\; lo \le w_i \le hi \quad \forall i \in 1\ldots N
\]

Los rangos $[lo, hi]$ varían según el conjunto de datos (véase Sección~\ref{sec:experimentos}).

%-----------------------------------------------------------------------
\section{Aplicación de los Algoritmos al Problema}
%-----------------------------------------------------------------------

\subsection{Representación de soluciones}

Una solución queda representada como un vector $w \in \mathbb{R}^N$
donde cada componente $w_i$ indica la fracción del capital invertida en el activo $i$.
La restricción de cardinalidad implica que $w_i$ puede ser exactamente cero
(no se invierte en esa empresa) o estar en el intervalo $[lo, hi]$.

\subsection{Generación aleatoria de soluciones}

Para construir una solución aleatoria válida se procede como sigue:

\begin{enumerate}
  \item Se elige aleatoriamente un número de empresas en las que \emph{no} se invierte,
        $n_0 \in [1, N{-}1]$.
  \item Los índices restantes $n_+ = N - n_0$ se seleccionan barajando un vector de índices.
  \item A cada empresa activa se le asigna un peso aleatorio en $[lo, hi]$.
  \item Se normaliza: $w_i \leftarrow w_i / \sum_k w_k$.
  \item Se llama al operador \fix{} para garantizar factibilidad.
\end{enumerate}

\subsection{Operador de reparación (\texttt{fix})}

Tras cualquier modificación de la solución puede violarse alguna restricción.
El operador de reparación actúa en tres fases:

\begin{enumerate}
  \item \textbf{Normalización:} $w_i \leftarrow w_i / \sum_k w_k$.
  \item \textbf{Recorte de rango:} si $0 < w_i < lo$ entonces $w_i \leftarrow 0$;
        si $w_i > hi$ entonces $w_i \leftarrow hi$.
  \item \textbf{Redistribución iterativa:} mientras $|\sum w_i - 1| > \varepsilon$:
        \begin{itemize}
          \item Si falta dinero por invertir: se reparte la diferencia entre las empresas
                activas con $w_i < hi$. Si no hay candidatos, se activa una empresa
                inactiva asignándole $lo$.
          \item Si sobra dinero: se retira la diferencia de las empresas con $w_i > lo$.
        \end{itemize}
\end{enumerate}

\subsection{Precómputo de matrices}

Para acelerar la evaluación, se precalculan antes de ejecutar los algoritmos:
el vector de factores de beneficio $b_i = \frac{1}{T}\sum_t \log(1+\mu_{ti})$
y la matriz de covarianza $\Sigma$, tanto para los datos de entrenamiento (2015--2024)
como para los de test (2025). Así, \texttt{fitness} se calcula en $O(N^2)$.

%-----------------------------------------------------------------------
\section{Algoritmos}
%-----------------------------------------------------------------------

\subsection{Algoritmo Aleatorio (\textit{Random Search})}

Genera \texttt{MAX\_EVALS} soluciones aleatorias independientes y devuelve la mejor.
Sirve como referencia inferior de calidad.

\begin{algorithm}[H]
\caption{Búsqueda Aleatoria}
\KwIn{problema, maxEvals}
\KwOut{mejor solución encontrada}
$\text{mejorFit} \leftarrow -\infty$\;
\Para{$i \leftarrow 1$ \KwTo $\text{maxEvals}$}{
  $s \leftarrow$ \createSolution{}\;
  $f \leftarrow$ \fitness{$s$}\;
  \Si{$f > \text{mejorFit}$}{
    $\text{mejorFit} \leftarrow f$\;
    $\text{mejorSol} \leftarrow s$\;
  }
}
\KwReturn{$\text{mejorSol}$}
\end{algorithm}

\subsection{Algoritmo Greedy}

Construye una solución de forma determinista asignando capital en orden decreciente
de una función heurística que estima el valor marginal de cada activo.

\subsubsection*{Función heurística}

La heurística para el activo $i$ estima su contribución neta al fitness si fuese
el único activo de la cartera, penalizando su covarianza media con el resto:
\[
  h(i) = b_i - \lambda \cdot \frac{1}{N}\sum_{j=1}^{N} \Sigma_{ij}
\]
donde $b_i$ es el factor de beneficio precalculado y $\Sigma_{ij}$ la covarianza.

\begin{algorithm}[H]
\caption{Greedy}
\KwIn{problema}
\KwOut{solución construida}
Calcular $h(i)$ para todo $i \in \{1,\ldots,N\}$\;
Ordenar empresas en orden decreciente por $h(i)$ $\rightarrow$ lista \textit{companies}\;
$w \leftarrow \mathbf{0}$;\quad $suma \leftarrow 0$\;
\Mientras{$suma < 1 - \varepsilon$ \textbf{y} quedan empresas en \textit{companies}}{
  $c \leftarrow$ primera empresa de \textit{companies}\;
  Eliminar $c$ de \textit{companies}\;
  $w[c] \leftarrow \min(hi,\; 1 - suma)$\;
  $suma \leftarrow suma + w[c]$\;
}
\KwReturn{$w$}
\end{algorithm}

\noindent El greedy es determinista: produce siempre la misma solución independientemente
de la semilla, por lo que su resultado es idéntico en las 50 ejecuciones.

\subsection{Búsqueda Local (BL)}

Implementa una búsqueda de \textbf{primero mejor} (\textit{first improvement})
sobre el vecindario definido por todos los pares de empresas $(p_1, p_2)$,
$p_1 \neq p_2$, con $p_1, p_2 \in \{1,\ldots,N\}$.

\subsubsection*{Operador de generación de vecino}

Dado un par $(origen, destino)$, el vecino se obtiene transfiriendo un
$ratio = 0{,}4$ del peso del activo origen al activo destino:
\[
  \text{change}(w, origen, destino) : \quad
  w'_{origen} = w_{origen} \cdot (1 - ratio),\quad
  w'_{destino} = w_{destino} + ratio \cdot w_{origen}
\]
Tras el cambio se aplica \fix{} para restaurar la factibilidad.

\subsubsection*{Esquema de la BL}

\begin{algorithm}[H]
\caption{Búsqueda Local (primero mejor)}
\KwIn{problema, maxEvals}
\KwOut{mejor solución encontrada}
$P \leftarrow$ todos los pares $(i,j)$ con $i \neq j$, $i,j \in \{1,\ldots,N\}$\;
$s \leftarrow$ \createSolution{};\quad $f \leftarrow$ \fitness{$s$};\quad $evals \leftarrow 1$\;
$mejora \leftarrow \text{verdadero}$\;
\Mientras{$mejora$ \textbf{y} $evals < maxEvals$}{
  $mejora \leftarrow \text{falso}$\;
  \shuffle{$P$}\;
  \Para{cada $(origen, destino) \in P$}{
    \Si{$evals \geq maxEvals$}{\textbf{romper}\;}
    $s' \leftarrow s$;\quad
    $s'[origen] \leftarrow s'[origen] \cdot 0{,}6$;\quad
    $s'[destino] \leftarrow s'[destino] + s[origen] \cdot 0{,}4$\;
    \fix{$s'$};\quad
    $f' \leftarrow$ \fitness{$s'$};\quad $evals \mathrel{+}= 1$\;
    \Si{$f' > f$}{
      $s \leftarrow s'$;\quad $f \leftarrow f'$;\quad $mejora \leftarrow \text{verdadero}$\;
      \textbf{romper}\;
    }
  }
}
\KwReturn{$s$}
\end{algorithm}

\noindent El vecindario tiene $N \cdot (N{-}1)$ vecinos. La exploración se baraja en
cada iteración para evitar sesgos de orden. La condición de parada es la convergencia
(ninguna mejora en un recorrido completo) o el agotamiento del presupuesto.

\subsection{Simulated Annealing (SA) --- Trabajo voluntario}

El Recocido Simulado extiende la búsqueda local aceptando soluciones peores con
probabilidad $e^{\Delta f / T}$, lo que permite escapar de óptimos locales.
La temperatura $T$ decrece geométricamente ($T \leftarrow T \cdot \alpha$),
lo que crea una transición natural de exploración (alta $T$) a explotación (baja $T$).

\subsubsection*{Calibración automática de temperatura}

En lugar de fijar $T_0$ manualmente, se estima automáticamente realizando 200
movimientos aleatorios, calculando el empeoramiento medio $\overline{\Delta}$
de los movimientos que empeoran, y fijando:
\[
  T_0 = \frac{-\overline{\Delta}}{\ln(0{,}8)}
\]
de modo que al inicio se acepta aproximadamente el 80\,\% de los movimientos
que empeoran, garantizando exploración amplia independientemente de la escala
del problema.

\subsubsection*{Operador de vecino}

Se elige aleatoriamente una empresa \textit{origen} con peso positivo y una empresa
\textit{destino} cualquiera distinta. Se transfiere una fracción aleatoria
$\delta \in [0{,}1, 0{,}9]$ del peso del origen al destino, seguida de \fix{}.
La aleatoriedad de $\delta$ proporciona mayor diversidad que el ratio fijo de la BL.

\begin{algorithm}[H]
\caption{Simulated Annealing}
\KwIn{problema, maxEvals, $\alpha = 0{,}999$}
\KwOut{mejor solución encontrada}
$T \leftarrow$ calibrar$T_0$(problema)\;
$s \leftarrow$ \createSolution{};\quad $f \leftarrow$ \fitness{$s$}\;
$s^* \leftarrow s$;\quad $f^* \leftarrow f$\;
\Para{$evals \leftarrow 1$ \KwTo $maxEvals - 1$}{
  $activas \leftarrow \{i : s[i] > 0\}$\;
  $origen \leftarrow$ aleatorio en $activas$\;
  $destino \leftarrow$ aleatorio en $\{1,\ldots,N\} \setminus \{origen\}$\;
  $s' \leftarrow s$;\quad
  $\delta \leftarrow$ uniforme$(0{,}1,\; 0{,}9)$\;
  $s'[origen] \mathrel{-}= \delta \cdot s[origen]$;\quad
  $s'[destino] \mathrel{+}= \delta \cdot s[origen]$\;
  \fix{$s'$};\quad $f' \leftarrow$ \fitness{$s'$}\;
  \Si{$f' > f$ \textbf{o} $\text{aleatorio}(0,1) < e^{(f'-f)/T}$}{
    $s \leftarrow s'$;\quad $f \leftarrow f'$\;
    \Si{$f > f^*$}{$s^* \leftarrow s$;\quad $f^* \leftarrow f$\;}
  }
  $T \leftarrow T \cdot \alpha$\;
}
\KwReturn{$s^*$}
\end{algorithm}

%-----------------------------------------------------------------------
\section{Estructura del Código y Manual de Usuario}
%-----------------------------------------------------------------------

\subsection{Estructura de directorios}

\begin{verbatim}
portfolio/
  main.cpp                  # Punto de entrada, experimentos de 50 semillas
  CMakeLists.txt            # Sistema de construcción
  boxplot.py                # Script Python para gráficas de cajas
  data/
    ibex_35.csv             # 30 empresas IBEX 35, 2015-2025
    syp_100.csv             # 97 empresas S&P 100, 2015-2025
    syp_500.csv             # 457 empresas S&P 500, 2015-2025
  common/
    mh.h                    # Interfaz genérica MH<tDomain>
    problem.h               # Interfaz genérica Problem<tDomain>
    solution.h              # Typedefs tSolution, tFitness
    random.hpp              # Generador pseudoaleatorio de alta calidad
  inc/
    portfolioProblem.h      # Declaración del problema
    randomSearch.h          # Algoritmo aleatorio
    greedy.h                # Algoritmo greedy
    localsearch.h           # Búsqueda local
    simulatedAnnealing.h    # Simulated Annealing (voluntario)
  src/
    portfolioProblem.cpp    # Implementación: carga datos, matrices, fitness
\end{verbatim}

\subsection{Compilación}

\begin{verbatim}
mkdir build && cd build
cmake ..
cmake --build .
\end{verbatim}

\subsection{Ejecución}

El ejecutable \texttt{main} no requiere argumentos. Ejecuta automáticamente los
cuatro algoritmos sobre los tres conjuntos de datos con 50 semillas cada uno,
imprime las tablas de medias por pantalla y genera los ficheros CSV para boxplot:

\begin{verbatim}
cd build
cmake --build . && cd ..
./build/main
\end{verbatim}

Al terminar se generan en el directorio raíz:
\texttt{results\_ibex\_35.csv}, \texttt{results\_syp\_100.csv},
\texttt{results\_syp\_500.csv}.

Para generar los boxplots (requiere Python con \texttt{pandas} y \texttt{seaborn}):

\begin{verbatim}
python3 boxplot.py
\end{verbatim}

Los parámetros configurables en \texttt{main.cpp} son \texttt{N\_SEEDS} (número
de semillas, por defecto 50) y \texttt{MAX\_EVALS} (presupuesto de evaluaciones,
por defecto 10\,000). Las semillas utilizadas son $1, 2, \ldots, 50$, iguales
para todos los algoritmos y conjuntos de datos.

%-----------------------------------------------------------------------
\section{Experimentos y Resultados}
\label{sec:experimentos}
%-----------------------------------------------------------------------

\subsection{Casos del problema}

\begin{center}
\begin{tabular}{lrrrl}
\toprule
\textbf{Fichero} & \textbf{Empresas ($N$)} & \textbf{Días train ($T$)} & \textbf{Días test} & \textbf{Rango $[lo, hi]$} \\
\midrule
\texttt{ibex\_35.csv}  & 30  & $\sim$2560 & 254 & $[0.005,\; 0.08]$ \\
\texttt{syp\_100.csv}  & 97  & $\sim$2516 & 249 & $[0.005,\; 0.05]$ \\
\texttt{syp\_500.csv}  & 457 & $\sim$2516 & 249 & $[0.005,\; 0.02]$ \\
\bottomrule
\end{tabular}
\end{center}

Parámetros comunes: $\lambda = 500$, \texttt{MAX\_EVALS} $= 10\,000$,
50 semillas ($1$--$50$), $\alpha_{\text{SA}} = 0{,}999$.

\subsection{Resultados}

Las tablas muestran la media de las 50 ejecuciones. El fitness reportado para el
Greedy es siempre el mismo al ser determinista.

\begin{table}[H]
\centering
\caption{Resultados IBEX 35 (medias, 50 semillas)}
\begin{tabular}{lrrrrr}
\toprule
\textbf{Algoritmo} & \textbf{Fit.\ 2015--24} & \textbf{Fit.\ 2025} & \textbf{Ben.\ 2025} & \textbf{Evals} & \textbf{Tiempo (s)} \\
\midrule
Greedy  & $-0.032$ & $-0.019$ & $-0.000$ & $1$         & $0.000$ \\
Random  & $-0.033$ & $-0.022$ & $0.000$  & $10\,000$   & $0.282$ \\
BL      & $-0.030$ & $-0.019$ & $-0.000$ & $10\,000$   & $0.259$ \\
SA      & $-0.030$ & $-0.019$ & $0.000$  & $10\,000$   & $0.187$ \\
\bottomrule
\end{tabular}
\end{table}

\begin{table}[H]
\centering
\caption{Resultados S\&P 100 (medias, 50 semillas)}
\begin{tabular}{lrrrrr}
\toprule
\textbf{Algoritmo} & \textbf{Fit.\ 2015--24} & \textbf{Fit.\ 2025} & \textbf{Ben.\ 2025} & \textbf{Evals} & \textbf{Tiempo (s)} \\
\midrule
Greedy  & $-0.025$ & $-0.031$ & $0.000$  & $1$         & $0.001$ \\
Random  & $-0.026$ & $-0.033$ & $0.000$  & $10\,000$   & $1.913$ \\
BL      & $-0.022$ & $-0.027$ & $0.000$  & $10\,000$   & $4.634$ \\
SA      & $-0.022$ & $-0.027$ & $0.000$  & $10\,000$   & $1.701$ \\
\bottomrule
\end{tabular}
\end{table}

\begin{table}[H]
\centering
\caption{Resultados S\&P 500 (medias, 50 semillas)}
\begin{tabular}{lrrrrr}
\toprule
\textbf{Algoritmo} & \textbf{Fit.\ 2015--24} & \textbf{Fit.\ 2025} & \textbf{Ben.\ 2025} & \textbf{Evals} & \textbf{Tiempo (s)} \\
\midrule
Greedy  & $-0.026$ & $-0.031$ & $0.000$  & $1$         & $0.038$ \\
Random  & $-0.028$ & $-0.039$ & $0.000$  & $10\,000$   & $42.966$ \\
BL      & $-0.019$ & $-0.028$ & $0.000$  & $10\,000$   & $199.447$ \\
SA      & $-0.019$ & $-0.027$ & $0.000$  & $10\,000$   & $35.846$ \\
\bottomrule
\end{tabular}
\end{table}

\subsection{Boxplots}

Las Figuras~\ref{fig:box_ibex}--\ref{fig:box_syp500} muestran la distribución del
fitness de entrenamiento en las 50 ejecuciones para cada conjunto de datos.

\begin{figure}[H]
  \centering
  \maybeimage{results_ibex_35.png}{width=0.7\textwidth}{results\_ibex\_35.png}
  \caption{Distribución del fitness (2015--2024) en IBEX 35.}
  \label{fig:box_ibex}
\end{figure}

\begin{figure}[H]
  \centering
  \maybeimage{results_syp_100.png}{width=0.7\textwidth}{results\_syp\_100.png}
  \caption{Distribución del fitness (2015--2024) en S\&P 100.}
  \label{fig:box_syp100}
\end{figure}

\begin{figure}[H]
  \centering
  \maybeimage{results_syp_500.png}{width=0.7\textwidth}{results\_syp\_500.png}
  \caption{Distribución del fitness (2015--2024) en S\&P 500.}
  \label{fig:box_syp500}
\end{figure}

%-----------------------------------------------------------------------
\section{Análisis de Resultados}
%-----------------------------------------------------------------------

\subsection{Calidad de las soluciones}

\textbf{BL supera claramente a Random y Greedy.} En ambos conjuntos con datos
completos (IBEX 35 y S\&P 100), la búsqueda local obtiene un fitness
aproximadamente un 11\,\% mejor que el algoritmo aleatorio y un 7\,\% mejor que
el greedy. Esto confirma que explotar la estructura local del espacio de búsqueda
es efectivo para este problema.

\textbf{El greedy ofrece una solución de calidad razonable con una única evaluación.}
Con solo 1 evaluación frente a las 10\,000 del resto, la heurística basada en
beneficio menos covarianza media supera al algoritmo aleatorio, lo que demuestra
que la función heurística captura información relevante del problema.

\textbf{SA y BL convergen a resultados prácticamente idénticos en IBEX 35 y S\&P 100.}
Esto se explica porque con $N$ pequeño (30 y 97 empresas) la BL ya es capaz de
recorrer el vecindario ($N(N{-}1)$ vecinos) en pocas iteraciones antes de agotar
el presupuesto, alcanzando óptimos locales de calidad comparable a los que
encuentra el SA.

\textbf{SA muestra mayor ventaja en S\&P 500 ($N=457$).} Con un vecindario de
457\,$\times$\,456\,$\approx$\,208\,000 pares, la BL apenas puede completar
un recorrido completo en 10\,000 evaluaciones y queda atrapada en el primer
óptimo local encontrado. El SA, al aceptar soluciones peores con probabilidad
controlada por la temperatura, puede escapar de esos óptimos y explorar regiones
más prometedoras del espacio.

\subsection{Variabilidad}

La BL presenta muy baja varianza entre semillas en los tres conjuntos (los
boxplots muestran cajas estrechas), lo que indica que el vecindario es suficiente
para guiar la búsqueda hacia zonas similares independientemente del punto de
partida. El SA hereda esta estabilidad y en S\&P 500 converge consistentemente
a valores de fitness similares entre semillas, algo que la BL no logra dado que
su resultado depende en mayor medida del punto inicial.

El algoritmo aleatorio presenta la mayor varianza, especialmente en S\&P 500
donde el espacio de búsqueda es enorme y 10\,000 muestras aleatorias cubren
una fracción ínfima de él.

\subsection{Coste computacional y escalabilidad}

\textbf{Coste por evaluación.}
La operación dominante es el cálculo de \texttt{fitness}, que requiere
$O(N^2)$ multiplicaciones (producto $w^\top \Sigma w$).
El operador \texttt{fix} añade un coste $O(N)$ por llamada, pero se vuelve
relevante cuando $N$ es grande y el proceso de redistribución de pesos necesita
más iteraciones para converger.

\textbf{Escalado empírico con $N$.}
La Tabla~\ref{tab:tiempos} recoge los tiempos medios por semilla de cada algoritmo.
Tomando S\&P 100 como referencia y comparando con S\&P 500
($N$ crece un factor $457/97 \approx 4{,}7$):

\begin{table}[H]
\centering
\caption{Tiempo medio por semilla (segundos) y factor de escala respecto a S\&P 100}
\label{tab:tiempos}
\begin{tabular}{lrrrrrr}
\toprule
 & \multicolumn{2}{c}{\textbf{IBEX 35} ($N=30$)}
 & \multicolumn{2}{c}{\textbf{S\&P 100} ($N=97$)}
 & \multicolumn{2}{c}{\textbf{S\&P 500} ($N=457$)} \\
\cmidrule(lr){2-3}\cmidrule(lr){4-5}\cmidrule(lr){6-7}
\textbf{Alg.} & t (s) & ×100 & t (s) & ref. & t (s) & ×100 \\
\midrule
Greedy & $0.000$ & ---    & $0.001$ & $1\times$  & $0.038$ & $38\times$ \\
Random & $0.006$ & $0.3\times$ & $0.038$ & $1\times$  & $0.859$ & $23\times$ \\
BL     & $0.005$ & $0.1\times$ & $0.093$ & $1\times$  & $3.989$ & $43\times$ \\
SA     & $0.004$ & $0.2\times$ & $0.034$ & $1\times$  & $0.717$ & $21\times$ \\
\bottomrule
\end{tabular}
\end{table}

\noindent El factor teórico esperado para un algoritmo cuyo coste por evaluación
es estrictamente $O(N^2)$ sería $(457/97)^2 \approx 22$.
Random y SA se aproximan a ese valor ($23\times$ y $21\times$),
confirmando que su coste está dominado por \texttt{fitness}.
La BL, en cambio, escala con un factor de $43\times$, sensiblemente mayor.
La razón es que la BL llama a \texttt{fix} en \emph{cada} vecino generado,
y con $N=457$ y \texttt{hi}=0{,}02 (máximo 50 activos simultáneos de 457),
la redistribución de pesos en \texttt{fix} requiere muchas más iteraciones
para alcanzar la factibilidad, elevando el coste efectivo por evaluación
por encima del teórico $O(N^2)$.

El Greedy escala con un factor $38\times$ de S\&P 100 a S\&P 500.
Aunque el greedy solo hace una evaluación, necesita precalcular
la heurística para cada empresa iterando sobre la fila correspondiente
de $\Sigma$ ($O(N^2)$ en total), lo que explica su tiempo no nulo en S\&P 500.

\textbf{Eficiencia calidad/tiempo.}
Una forma de comparar algoritmos con distinto presupuesto temporal es
la mejora de fitness obtenida por unidad de tiempo.
En S\&P 500, la BL tarda 199{,}4\,s por semilla y alcanza un fitness de $-0{,}019$,
mientras que el SA obtiene el mismo fitness en solo 35{,}8\,s, es decir,
con \textbf{un coste 5{,}6 veces menor}.
Esto se debe a la diferente estrategia de generación de vecinos:
la BL construye y baraja los $N(N{-}1) \approx 208\,000$ pares en cada
iteración ($O(N^2)$ de overhead estructural), mientras que el SA
selecciona aleatoriamente un par $(origen, destino)$ en $O(1)$,
utilizando todo el presupuesto en evaluaciones de fitness en lugar de en
gestión del vecindario.

\textbf{Conclusión sobre eficiencia.}
Para instancias pequeñas ($N \lesssim 100$) los cuatro algoritmos son
prácticamente inmediatos y BL y SA ofrecen idéntica calidad.
Para instancias grandes ($N \sim 500$), el SA es la opción más eficiente:
iguala la calidad de la BL consumiendo un quinto del tiempo,
gracias a su generación de vecinos en $O(1)$ frente al $O(N^2)$
de estructurar el vecindario completo de la BL.

%-----------------------------------------------------------------------
\begin{thebibliography}{9}
\bibitem{markowitz1952} H.~Markowitz, ``Portfolio selection,''
  \textit{Journal of Finance}, vol.~7, no.~1, pp.~77--91, 1952.

\bibitem{kirkpatrick1983} S.~Kirkpatrick, C.~D.~Gelatt, and M.~P.~Vecchi,
  ``Optimization by simulated annealing,''
  \textit{Science}, vol.~220, no.~4598, pp.~671--680, 1983.

\bibitem{random_hpp} E.~Konkov, \textit{random.hpp: A C++11/14/17 header-only
  random number generator library}, GitHub, 2018.
\end{thebibliography}

\end{document}
